\documentclass[11pt,a4paper]{article}
\usepackage[margin=1in]{geometry}
\usepackage{amsmath,amssymb,amsthm,mathtools}
\usepackage{enumitem}
\usepackage{hyperref}
\usepackage{xcolor}
\usepackage{booktabs}
\usepackage{array}
\usepackage{longtable}
\usepackage{listings}
\usepackage{titlesec}
\usepackage{microtype}
\newcommand{\MG}{\textsc{MathGraph}}
\newcommand{\LOGOS}{\textsc{LOGOS}}
\newcommand{\Chora}{\textsc{Chora}}
\newcommand{\Lawbook}{\textsc{Lawbook}}
\newcommand{\ObstructionAtlas}{\textsc{Obstruction Atlas}}
\newcommand{\ReasonAtlas}{\textsc{Reason Atlas}}
\newcommand{\RootAtlas}{\textsc{Root Node Atlas}}
\newcommand{\HTilt}{\textsc{H-Tilt}}
\newcommand{\Lean}{\textsc{Lean}}
\newcommand{\Isabelle}{\textsc{Isabelle}}
\newcommand{\ETP}{\textsc{ETP}}
\newcommand{\SAIR}{\textsc{SAIR}}
\title{\textbf{MathGraph}\\A Full Vision, Design, and Technical Specification\\\large A Generative Verification Kernel for Lawful Mathematical Continuation}
\author{Heath Sanchez}
\date{}
\begin{document}
\maketitle
\begin{abstract}
\MG{} is a lightweight generative verification kernel for lawful mathematical continuation. It separates proposal from the verifier boundary and accepts claims only through a proof, finite countermodel, or named obstruction terminal contract. Around that boundary it builds proof digestion, process memory, \Lawbook{} acceptance and query, typed projection, role-based object introduction, and structural analogy with exposition. These layers improve reuse, explanation, and compounding discovery while remaining advisory unless verifier, trusted-importer, finite-validation, or chain-audit evidence promotes truth.
\end{abstract}

\section{One-Line Vision}
\MG{} is a lightweight generative verification kernel where models propose, \MG{} constrains, verifiers decide, digestion assimilates, the \Lawbook{} remembers, and projection scales.

\noindent LLMs dream. \Chora{} receives. Agents explore. \HTilt{} selects. Alchemy transforms. \LOGOS{} constrains. Verifiers decide. The \Lawbook{} remembers. Projection scales.

\section{Core Doctrine}
\begin{description}[style=nextline]
\item[Continuation Under Constraint] Search may be generative, but continuation is lawful only under explicit typed constraints.
\item[Verifier Boundary] Only verifier, trusted importer, finite validator, and chain audit boundaries promote truth.
\item[Terminal Contract] Accepted claims end in exactly one terminal form.
\item[Compounding] Verified work should become reusable memory that improves future work.
\item[Advisory / Truth Separation] All non-verifier artifacts are advisory: a thing can guide search without being true, explain without verifying, schedule without proving, and be accepted as advisory memory without becoming a theorem.
\end{description}

\section{Why MathGraph Exists Now}
AI proof generation is becoming abundant. Verification is necessary but insufficient; proof digestion is the missing third stage. Mathematical workflows need infrastructure, not merely faster theorem generation. Correct but indigestible proofs can fail the system-level goal. The relevant metric is not merely ``proved,'' but ``verified, digested, reusable, explainable, and able to improve future work.'' \MG{} is infrastructure for the proof-abundance era.

\section{Formal Terrain}
The universal proof hypergraph contains possible objects, claims, transformations, witnesses, and derivations. A structural hypergraph records reusable shapes across that space. Human mathematics is a small ribbon through proof space, selected by importance rather than raw reachability. Importance combines compression, reuse, projection gain, digestion value, and explanation value.

\section{Terminal Forms and Trust Contract}
\begin{itemize}
\item \texttt{VERIFIED\_PROOF}: proof artifact accepted by a formal boundary.
\item \texttt{FINITE\_COUNTERMODEL}: checked finite witness preserving source and separating target.
\item \texttt{NAMED\_OBSTRUCTION}: structured accepted obstruction memory.
\item residual/advisory status: useful history or pressure that has not crossed truth boundaries.
\end{itemize}

\section{Kernel Architecture}
\begin{lstlisting}
Domain Claim
  -> Formal-World Adapter Registry
  -> Adapter Capability / Parse / Normalize / Validate
  -> Continuation Actions
  -> Continuation Curriculum
  -> Verification Episode
  -> Verifier / Importer / Finite Validator / Chain Audit
  -> Proof Digestion
  -> Verifier Feedback / Repair
  -> Discovery Value
  -> Lawbook Acceptance
  -> Lawbook Query / Known Skip
  -> Structural Identity
  -> Habit Rules
  -> Reason Compression
  -> Process Memory
  -> Structure Registry / Typed Projection
  -> Role-Based Object Introduction
  -> Structural Analogy / Exposition
  -> Projection
  -> Telemetry
\end{lstlisting}

\section{Alchemical Process}
Raw Matter, Calcination, Solution, Sublimation, Descension, Distillation, Coagulation, Fixation, Ceration, Conjunction, Multiplication, Projection, and Perfection name process stages. Alchemical traces are process records. They may record fixation, but only verifier, trusted importer, finite validator, or chain audit evidence promotes truth.

\section{H-Tilt and Route Pressure}
Soft tilt ranks work heuristically. Killed-dynamics or spectral \HTilt{} estimates survival from route telemetry. \HTilt{} Lite uses taste, scars, cost, compression, and projection gain. Discovery value adds route-level utility signals. All are scheduling pressure, not truth.

\section{Agentic Layer}
\texttt{AgentProfile} and \texttt{AgentExperience} record lightweight agents, taste, scars, and biography. Current agents are advisory policies, not truth authorities. Existential agents remain future work.

\section{Lawbook and Accepted Memory}
The \Lawbook{} is the accepted public-memory boundary. Candidate entries differ from accepted entries. Accepted truth entries require boundary evidence; digestion, projection, and discovery value may recommend entries but cannot turn themselves into truth.

\section{Lawbook Query and Known-Skip Service}
Accepted verified proofs can skip, accepted finite countermodels can skip, and accepted obstructions can skip only as accepted obstruction memory. Candidate-only, digestion-only, projection-only, analogy-only, and role-only records cannot skip as truth.

\section{Proof Digestion}
The workflow is proof generation, proof verification, then proof digestion. Digestion produces dependency maps, step classification, key idea candidates, reusable schema candidates, exposition notes, and projection hints. It inherits verified boundaries; it does not create them.

\section{Verifier Feedback and Repair}
Feedback records flaw severity. Repair plans may locally revise, reroute, emit proof tasks, emit obstruction tasks, hold in \Chora{}, or residualize. Feedback and repair plans do not verify claims.

\section{Continuation Actions and Curricula}
Hard targets decompose into warm-ups, simplified cases, finite examples, proof tasks, countermodel tasks, projection tasks, repair tasks, and replayable episode inputs. Curricula are planning artifacts only.

\section{Discovery Value}
Discovery value scores curricula, digestion, feedback, repair, projection, actions, alchemy, agents, and telemetry. Value is route pressure, not proof.

\section{Structural Identity and Canonicalization}
Typed structural graph views, deterministic signatures, duplicate detection, and advisory merge/conflict candidates keep memory tidy. Signature matches or isomorphism-looking matches are not equality or proof.

\section{Habit Rules and Route Promotion}
Repeated successful route outcomes become habit candidates. Reviewed habits adjust scheduling pressure advisably; habits do not verify claims.

\section{Reason Compression}
Reason compression proposes minimal sufficient reasons, distinguishes load-bearing atoms from decorative atoms, and emits reason nodes. Reasons are explanatory pressure, not theorem proofs.

\section{Process Memory}
Process episodes retain included context, excluded context, transitions, eliminations, inherited boundaries, and route/process summaries. Process memory records how work unfolded; it does not verify.

\section{Structure Registry and Typed Projection}
Broad families include algebraic, order, topological, combinatorial, logical, computational, graphical, dynamical, probabilistic, categorical, and mixed structures. Compatibility may be exact or same-family, require cross-family review, need an adapter, need formalization, block on type mismatch, or block on conflict. Typed projection is advisory.

\section{Role-Based Object Introduction}
Recurring role signatures yield definition candidates, witness candidates, conjecture/task candidates, review, and accepted advisory role objects. Role objects are concepts, not proofs or certificates.

\section{Structural Analogy and Exposition}
Analogy sources are compared through feature maps. Breaks, candidates, exposition notes, and review create projection, digestion, formalization, adapter, and review tasks. Analogy and exposition are understanding aids, not verification.

\section{Richer Formal-World Adapters}
M7 adds a lightweight typed gate between advisory metabolism and concrete
formal systems. Formal-world adapters detect broad world kinds, expose adapter
specifications and capability reports, parse raw inputs, normalize
representations, validate formal shape, and generate proof, countermodel,
formalization, finite-validation, adapter, and review tasks. They prepare
explicit handoffs to a verifier, trusted importer from verified external
artifacts, finite validator, or chain audit. Adapter-aware typed projection,
role, and analogy tasks can now state when an adapter, formalization step, or
external boundary is required. Parse, normalize, and validate success remains
advisory; an adapter does not verify unless an external boundary result is
returned with explicit evidence.

\section{Derived Closure and Projection}
Derived closure uses transitivity, source weakening, target strengthening, equivalence collapse, duality where justified, and constructor replay. FALSE-side certificate directionality must be preserved: projection may reuse verified memory, never reverse an invalid implication into truth.

\section{SAIR / ETP Nursery World}
\ETP{} over magmas is the nursery and \SAIR{} proving ground. The FALSE side uses finite countermodels. The TRUE side requires proof verification. Finite-search failure is not TRUE. \SAIR{} is the first proving ground, not the whole product.

\section{API Surfaces}
Current backend-first surfaces include lawbook query, verification episode, projection, roadmap alignment, proof digestion, verifier feedback, continuation curriculum, discovery value, structural identity, habit rules, reason compression, process memory, structure registry, role objects, and structural analogy.

\section{Roadmap Alignment Checker}
Criticals include advisory output claiming \texttt{VERIFIED\_PROOF} or \texttt{FINITE\_COUNTERMODEL}; finite-search miss represented as proof; route score treated as truth; digestion treated as proof; feedback text treated as boundary; curriculum warm-up treated as target proof; discovery value treated as truth; structural signature treated as equality; habit treated as verifier; reason node treated as theorem; process memory treated as verifier; typed projection treated as implication; role object treated as theorem; analogy or exposition treated as verification; and terminal form without boundary evidence. Positive signals preserve advisory boundaries and record lawful reuse.

\section{Metrics}
\begin{longtable}{ll}
\toprule
Metric & Purpose\\
\midrule
CertificateYieldPerCompute & terminal yield per cost\\
ResidualCompressionGain & residual reduction\\
DerivedAmplificationFactor & lawful reuse expansion\\
ProjectionGain & verified memory reuse\\
KnownSkipRate & accepted-memory avoidance\\
RoutePrecision & useful route selection\\
ObstructionNamingRate & structured failure capture\\
ConstructorReuseRate & constructor compounding\\
ProofDigestionYield & intelligibility gain\\
RepairLoopYield & repair productivity\\
CurriculumStageYield & staged task productivity\\
DiscoveryValueCalibration & value-score reliability\\
StructuralDuplicateAvoidance & memory hygiene\\
HabitPrecision & habit usefulness\\
ReasonCompressionGain & explanation compression\\
ProcessReplayCoverage & replayable history\\
TypedProjectionSafety & unsafe transfer avoidance\\
RoleReuseRate & concept reuse\\
AnalogyExpositionYield & readable relation yield\\
LawbookCompression & accepted-memory density\\
\bottomrule
\end{longtable}

\section{Efficiency Discipline}
Documentation must not overclaim code. Every advisory layer must be deletable without breaking truth. Expensive frontier models are not the first line of defense: use cheap deterministic filters first and formal verifiers where boundary crossing is required. Keep public terminology clean and remain API/backend first.

\section{Implementation Roadmap}
\begin{longtable}{ll}
\toprule
Milestone & Scope\\
\midrule
M0 & Closed-Loop Certificate Smoke\\
M1 & Agentic Alchemical Loop\\
M2 & Projection Engine\\
M3 & Root-Aware Constructors\\
M4 & TRUE-Side Proof Verification\\
M4.5 & Unified Verification Episode\\
M5 & Route Telemetry and Spectral H-Tilt\\
M6 & Domain-General Claim IR\\
M6.5 & Lean Adapter Hardening\\
M6.6 & Continuation Actions\\
M6.7 & Proof Digestion\\
M6.8 & Verifier Feedback / Repair\\
M6.9 & Continuation Curriculum\\
M6.10 & Discovery Value\\
M6.11 & Lawbook Acceptance Boundary\\
M6.12 & Lawbook Query / Known Skip\\
M6.13 & Structural Identity\\
M6.14 & Habit Rules\\
M6.15 & Reason Compression\\
M6.16 & Process Memory\\
M6.17 & Structure Registry / Typed Projection\\
M6.18 & Role-Based Object Introduction\\
M6.19 & Structural Analogy / Exposition\\
M6.20 & Documentation / Spec Synchronization and Publication Refresh\\
M7 & Richer Formal-World Adapters\\
\midrule
M8 & Deeper Proof-System Integration\\
M9 & Semantic and Natural-Language Claim Handling with Strict Boundaries\\
M10 & API Service Hardening\\
M11 & Existential Agent Ecology\\
\bottomrule
\end{longtable}

\section{Risks and Guardrails}
Guard against proof generation without digestion; verification without understanding; advisory pressure mistaken for proof; route scores treated as truth; finite-search failure mistaken for TRUE; natural-language feedback treated as verifier boundary; proof digestion treated as proof; warm-up curricula treated as target proof; discovery value treated as truth; structural identity treated as proof equality; habits treated as authority; reasons treated as theorems; process memory treated as verification; typed projection treated as implication; role objects treated as definitions or theorems; analogy treated as proof; exposition treated as verification; stale docs causing public overclaim; and external branded terms confusing \MG{}'s identity.

\section{Final Compression}
\MG{} is the world. Agents are the explorers. \HTilt{} is the compass. Alchemy is the method. \LOGOS{} is the law. Verification is the tribunal. Digestion is understanding. Process memory is biography. The \Lawbook{} is memory. Projection is scaling. Analogy is exposition, not truth.

\MG{} should not scale by becoming bigger. It should scale by making every verified thing reusable and every failure informative.
\end{document}
